\section{Discussion}\label{sec:discussion}

\subsection{What kind of learner this turned out to be}\label{sec:disc-kind}

v3 is not one mechanism but three learners under one convex stacker, and the input type decides which of them carries the fit (\cref{sec:learner}): enumerate-and-test on Boolean inputs, at about an Occam sample cost wherever its search can propose the right feature \status{Heuristic}; a fixed, data-derived hierarchy of subset sums read out by gradient-boosted trees on images, with nothing fitted inside the features; and the classical stack of Part 8 on other tabular data. That is the AutoML-style dispatch by input type of the Parts 6--9 paper (\cref{sec:relates}), and Part 10's two additions, a wider sum-block search for DFG and a second booster over a richer bank, sit inside the first two regimes. What the results add is where each regime's performance comes from. On Boolean inputs, the interpretability analysis places DFG's data cost in its search's low-degree seeds rather than in its test (\cref{sec:interp-toy}; a reading of descriptive measurements, not an ablation). On images, the gain over v2 is a better fixed representation, which ties a CPU-matched MLP ensemble at $n=60000$ (\cref{sec:res-img-gain}). On tabular data nothing has changed since Part 8 (\cref{sec:res-tabular}).

\subsection{Why it works where it works}\label{sec:disc-why}

\paragraph{Boolean targets: an Occam-like sample cost, when the search can propose} \status{Heuristic}. DFG's test admits a feature once $n$ exceeds a multiple of the log-size of its class over the squared residual correlation (\cref{sec:dfg}). For a $k$-sparse character, $\ln N\approx k\ln d$; for a signed block on $k$ of $d$ inputs, $\ln N\approx\ln\bigl(\binom dk2^{k-1}\bigr)$. Either way the sample cost is logarithmic in the size of the dictionary: the Occam rate, paid for with FLOPs of enumeration and search that fit inside one network run at $d=30$. Gradient-trained networks pay in samples instead. For parities the companion paper measured rates close to the leap and correlational-SQ rate $d^{k-1}$ \cite{gdcircuit2026,abbe2022msp,abbe2023leap,barak2022hidden,kearns1998sq}. For an even function of one signed sum, a network has to find the block's sign vector in the degree-2 statistic $\E[f\,xx^\top]$ against sampling noise. That is the information-exponent picture of single-index learning \cite{benarous2021online,damian2024gaps}. Networks can trade data for width and compute along a frontier \cite{edelman2023pareto}; DFG sits at an extreme of that trade.

The interpretability analysis sharpens the picture in one respect (\cref{sec:interp-toy}). On the pure symmetric targets the test is not the bottleneck, and the data cost is consistent with being set by when the search first \emph{proposes} the block, which it seeds from degree-1 and degree-2 statistics of the residual (findings 2 and 3 of \cref{sec:interp}). So DFG relies on the same low-order statistics that a network's gradient sees. It needs far fewer samples because it rounds a coarse spectral estimate to an exact signed subset and then tests that subset with a union bound, where a network has to learn the direction to useful precision by noisy local steps \status{Heuristic}.

\paragraph{Where neither has a foothold.} The residue classes mod 3 of 10--14-input counts have almost no Fourier weight at degrees 1 and 2, and DFG's search never proposes their block, although the test would accept it long before $n=16384$ (\cref{sec:interp-toy-fail}). The networks mostly fail too: none learns f7\_05, and only the most regularized variant learns h6\_00, at $\nstar=12304$ (fd\_10, in \code{dlfamdev}, has no network runs). Their signal sits at high degree on an unknown subset of inputs, which learners that rely on low-order correlations should find hard \cite{kearns1998sq}. Exhaustive search does not help within the budget: there are $1.5\times10^{10}$ to $1.2\times10^{12}$ blocks of the true size, each costing at least $n$ operations to score, against $1.62\times10^{11}$ FLOPs per toy fit. So this is a shared limit, not a loss specific to discrete search \status{Heuristic}. The only target on which a network is clearly better than v3 is of this kind.

\paragraph{Images: locality recovered from data, read out by trees} \status{Heuristic}. On QuickDraw the most correlated pixels of a pixel are its neighbours, so the pooling sets are small blobs of the true grid and the bars short runs along it (\cref{sec:interp-img-geom}). Recovering image topology from such correlations is an old idea \cite{leroux2007topology}, as is grouping inputs into receptive fields by correlation \cite{coates2011receptive}. Boosted trees \cite{ke2017lightgbm} over these fixed local sums then do what a first convolutional layer with a nonlinearity would do, but at fixed positions and without shared weights. The MLP references are denied this structure by design: they are invariant to the pixel permutation, and nothing in their training recovers it. So v3 beats MLP ensembles that spend 2.0 and 5.0 times its FLOP budget (\code{mlp_strong} at all six held-out cells, \code{mlp_wide} at five of six), and ties the 1024-wide ensemble, which spends 9.2 budgets and a comparable time per fit. The control and context runs bound how much of this is representation (\cref{sec:res-img-gain}): an MLP given v2's pooled sums ties v2 but not v3, no MLP was given the bars, and a CNN told the layout (seed 0) has lower error at all three cells where both ran, clearly at one. It does not show that discrete search beats gradient-based learning on images.

\paragraph{Tabular data} \status{Heuristic}. Here the result is the known strength of tree ensembles and stacking on medium-sized tabular data against standard-recipe MLPs (\cref{sec:res-tabular}); DFG runs on 2 of the 26 datasets and changes no cell there clearly.

\subsection{Limits}\label{sec:disc-limits}

\begin{itemize}
\item \textbf{Small held-out sets.} \code{dlheld7} has 12 targets from one family, and on 5 of its 10 wins default gradient boosting or logistic regression already beats the networks. \code{qd_c} and \code{qd_d} are two 10-class sets from one source. The 10 fresh tabular datasets of Part 8 give $p=0.11$. A second family or source could change the picture.
\item \textbf{The earlier suites are in-sample} (\cref{sec:discipline}). Pooled numbers over all suites summarize the record; they are not held-out estimates. The margin against \code{nets_wd} is smallest on the two draws of the blind family, which is what selection on the other suites would produce, though the suites also differ in kind and the differences are within noise (\cref{sec:res-toy-all}).
\item \textbf{One seed at $n\ge30000$ on images.} The two cells that decide the comparison with \code{mlp_xl} rest on one fit each, and on the test set's binomial noise for their standard error. v3's gain over v2 is below the 10\% threshold at every held-out cell; that v3 is lower in all 14 seed-paired fits is a description, not a pre-registered test (\cref{sec:res-img-heldout}).
\item \textbf{Targets without a low-degree foothold} are found late or never: the residue classes mod 3 never, h5s7 at $n=8192$, and h5s1 in no learner's 3-seed mean (\cref{sec:res-toy-fail,sec:interp-toy-fail}). No cheap search for them is known (\cref{sec:disc-next}).
\item \textbf{Locality is used in spirit.} The image members never read a pixel index, but the correlations give back neighbourhoods and lines, the information the permuted-pixel protocol was meant to withhold (\cref{sec:interp-img-geom}). No MLP was given the bars, and a CNN told the layout has lower error wherever both ran (\cref{sec:res-img-gain}).
\item \textbf{Equal FLOPs are not equal CPU time, and prediction is not charged.} v3 uses 0.30--0.58 of its FLOP budget on the held-out cells, but takes about as long per fit as \code{mlp_xl}, which it only ties. At $n=3000$ it is close to the budget guard once Part 9's audited correction is applied, and no compute audit of v3 was made. A v2 query costs about $500\times$ an MLP's, and v3 adds to it (\cref{sec:res-img-compute}).
\item \textbf{Concrete references.} One toy network architecture at one width and one learning rate, and image MLPs without augmentation or architecture search. The hindsight minima favour the references, but stronger networks would plausibly narrow the margins.
\item \textbf{No independent audit in Part 10.} Parts 8 and 9 had independent audits of leakage, compute, discreteness and reference strength \cite{discretelearners2026}. Part 10 inherits their findings and fixes, but its own comparisons were not re-audited in that sense. The interpretability analyses were independently verified, and the paper's claims were checked against the result files and records by independent checkers (\cref{app:provenance}).
\item \textbf{The interpretability is descriptive}: seed 0, no ablation of the learner or causal test, mechanism labels for v1.3's changes read off fitted models, and stacking weights, split gain and contribution maps that describe the fits, not what their predictions depend on (\cref{sec:interp}).
\end{itemize}

\subsection{What would test it further}\label{sec:disc-next}

\begin{enumerate}
\item \emph{The information-matched image comparison.} Give an MLP v3's whole bank, pooled sums and bars, as \code{mlp_pool} did for v2's. Train a gradient-based network that learns sparse local connectivity \cite{neyshabur2020convolutions}, or a locally connected network on the recovered neighbourhoods. Either would say how much of v3's image margin is representation that a network can use as well.
\item \emph{A second blind toy family}, written independently of the first, and the toy suites at larger $d$. At $d=30$ the advantage on the blind family is a factor of about $2^{1.5}\approx2.8$. For parities DFG's advantage grew with $d$ in Part 6; whether it does for signed sums was not measured.
\item \emph{Block search beyond degree 2.} A seed that reads higher-order statistics of the residual could reach the mod-$q$ targets. The Part 10 builder found that the true block of fd\_10 is detectable only within about 3--4 errors of the truth, so that random restarts would need far more than $10^5$ tries; a cheaper route is not known.
\item \emph{Power and attribution.} More seeds at $n=60000$; a pre-registered paired comparison with v2; image sets from another source; ablations of v1.3's three additions on held-out targets; and a compute audit of v3 that also charges prediction.
\end{enumerate}
